%%%%%%%%%%%%%%%%%%%%%%%%%%%%%%%%%%%%%%%%%%%%%%%%%%%%%%%%%%%%%%%%%%%%%%
% Lớp 10 | Phương pháp tọa độ trong mặt phẳng (𝑂𝑥𝑦) | Đường thẳng trong mặt phẳng tọa độ
% Dạng 0H9?3-4: Bài toán về góc giữa hai đường thẳng
% Số câu: 19  (sinh tự động bởi tach_cau.py; từ nay thầy sửa trực tiếp file này)
%%%%%%%%%%%%%%%%%%%%%%%%%%%%%%%%%%%%%%%%%%%%%%%%%%%%%%%%%%%%%%%%%%%%%%

%%%%%%%%%%%%%%%% TRẮC NGHIỆM 4 PHƯƠNG ÁN %%%%%%%%%%%%%%%%

% nguồn: MaHoaNH_TracNghiem_Lop10
\begin{ex}%[0H9N3-4]%[ma:f1635ac6]
	Trong mặt phẳng tọa độ $Oxy$, góc  giữa hai đường thẳng $d_1 \colon 3x+y+3=0$ và $d_2 \colon -2x+y+1=0$ là
	\choice 
		 {$135^\circ$}
		 {\True $45^\circ$}
		{$90^\circ$} 
		 {$30^\circ$}
	 \loigiai{
		Ta có $\cos (d_1,d_2)=\dfrac{|3 \cdot  (-2)+1\cdot 1|}{\sqrt{3^2+1^2}\cdot \sqrt{(-2)^2+1^2}}=\dfrac{\sqrt{2}}{2}$.\\
		Suy ra góc giữa hai đường thẳng $d_1$ và $d_2$ là  $45^\circ$.
	} 
\end{ex}

% nguồn: MaHoaNH_TracNghiem_Lop10
\begin{ex}%[0H9N3-4]%[ma:025324f0]
	Trong mặt phẳng $Oxy$,  góc giữa hai đường thẳng $d_1\colon 2x-y-10=0$ và $d_2\colon x-3y+9=0$ bằng
	 \choice
		 { $30^\circ$}
		{$60^\circ$} 
		 {\True $45^\circ$}
		{$135^\circ$} 
	\loigiai{ Đường thẳng $d_1\colon 2x-y-10=0$ có véc-tơ pháp tuyến  $\overrightarrow{n}_1=(2;-1)$.\\
		 Đường thẳng $d_2\colon x-3y+9=0$ có véc-tơ pháp tuyến $\overrightarrow{n}_2=(1;-3)$.\\
		 Ta có $\cos \left(d_1, d_2\right)=\dfrac{|2\cdot 1+(-1)(-3)|}{\sqrt{2^2+(-1)^2}\cdot \sqrt{1^2+(-3)^2}}=\dfrac{\sqrt{2}}{2}$.\\
		Góc giữa hai đường thẳng $d_1$  và $d_2$ bằng $45^\circ$.
	 }
\end{ex}

% nguồn: MaHoaNH_TracNghiem_Lop10
\begin{ex}%[0H9N3-4]%[ma:b5665cb6]
	Trong mặt phẳng tọa độ $Oxy$,  tính góc giữa hai đường thẳng $\Delta\colon x-\sqrt{3}y+2=0$ và $\Delta'\colon x+\sqrt{3}y-1=0$.
	\choice 
		 {\True $60^{\circ}$}
		 {$30^{\circ}$}
		{$90^{\circ}$} 
		{$120^{\circ}$} 
	 \loigiai{
		Ta  có
		\[ 
		 \cos(\Delta,\Delta')=\dfrac{|1\cdot 1+(-\sqrt{3}\cdot \sqrt{3})|}{\sqrt{1^2+(-\sqrt{3})^2}\cdot\sqrt{1^2+(\sqrt{3})^2}}=\dfrac{1}{2}\Rightarrow (\Delta,\Delta')=60^{\circ}.
		\] 
	 }
\end{ex}

% nguồn: MaHoaNH_TracNghiem_Lop10
\begin{ex}%[0H9H3-4]%[ma:6e293907]
	Cho  hai đường thẳng $d_1 \colon \heva{&x=1-2t_1 \\ &y=2+t_1}$ và $d_2 \colon \heva{&x=2+t_2 \\ &y=5+2t_2}$. Số đo góc giữa hai đường thẳng $d_1$ và $d_2$ bằng
	\choice 
		 {$45^{\circ}$}
		 {$60^{\circ}$}
		 {\True $90^{\circ}$}
		{$135^{\circ}$} 
	 \loigiai{
		Vectơ chỉ phương của đường thẳng $d_1$, $d_2$ lần lượt là $\vec{u}_1=(-2;1)$, $\vec{u}_2=(1;2)$. \\ 
		Ta có $\vec{u}_1 \cdot \vec{u}_2 = (-2) \cdot 1 + 1 \cdot 2 = 0  \Rightarrow d_1 \perp d_2$.
	} 
\end{ex}

% nguồn: MaHoaNH_TracNghiem_Lop10
\begin{ex}%[0H9H3-4]%[ma:23f09c7a]
	Góc giữa hai đường  thẳng $\Delta_{1}\colon3x-2y+1=0$; $\Delta_{2}\colon\heva{&x=t\\&y=7-5t}$ $(t\in \mathbb{R})$.
	 \choice
		{\True $45^\circ $} 
		{$30^\circ  $}
		{$135^\circ $} 
		 {$60^\circ $}
	 \loigiai{
		$\Delta_{1}$ có  $1$ véc-tơ pháp tuyến là $\overrightarrow{n}_{1}=\left(3;-2\right)$.\\
		$\Delta_{2}$ có $1$ véc-tơ chỉ phương là $\overrightarrow{u}=\left(1;-5\right)$ nên có $1$ véc-tơ pháp tuyến là $\overrightarrow{n}_{2}=\left(5;1\right)$.\\ 
		Do đó  $\cos(\Delta_{1},\Delta_{2})=\dfrac{\left|\overrightarrow{n}_{1}\cdot\overrightarrow{n}_{2}\right|}{\left|\overrightarrow{n}_{1}\right|\cdot\left|\overrightarrow{n}_{2}\right|}=\dfrac{\sqrt{2}}{2} \Rightarrow \left(\Delta_{1},\Delta_{2}\right)=45^\circ$.
	} 
\end{ex}

% nguồn: MaHoaNH_TracNghiem_Lop10
\begin{ex}%[0H9H3-4]%[ma:6bd325e1]
	Tìm góc giữa hai đường thẳng $\Delta_{1}\colon x-2y-5=0$ và $\Delta_{2}\colon \heva{&x=2-t\\&y=3+2t}\quad(t\in \mathbb{R}).$ 
	 \choice
		{\True $90^{\circ}$} 
		{$60^{\circ}$} 
		{$45^{\circ}$} 
		{$30^{\circ}$} 
	\loigiai{ 
		Đường thẳng $\Delta_{1}\colon  x-2y-5=0$ có vectơ pháp tuyến là $\overrightarrow{n}_{1}=(1;-2)$.\\
		Đường thẳng $\Delta_{2}\colon \heva{&x=2-t\\&y=3+2t}\quad(t\in \mathbb{R})$ có vectơ chỉ phương là $\overrightarrow{u}_{2}=(-1;2)\Rightarrow$ vectơ pháp tuyến là $\overrightarrow{n}_{2}=(2;1)$.\\ 
		Ta có $\cos \left(\Delta_{1},\Delta_{2}\right)=\dfrac{\left|\overrightarrow{n}_{1}\cdot \overrightarrow{n}_{2}\right|}{\left|\overrightarrow{n}_{1}\right|\cdot\left|\overrightarrow{n}_{2}\right|}=\dfrac{|1\cdot 2+(-2)\cdot 1|}{\sqrt{1^{2}+(-2)^{2}}\cdot  \sqrt{2^{2}+1^{2}}}=0$.\\
		Suy  ra $\left(\Delta_{1},\Delta_{2}\right)=90^{\circ}$.
	} 
\end{ex}

% nguồn: MaHoaNH_TracNghiem_Lop10
\begin{ex}%[0H9H3-4]%[ma:0f4d0933]
	Tính góc giữa  hai đường thẳng $d_{1}\colon 3x+y-1=0$ và $d_{2}\colon 4x-2y-4=0$.
	\choice 
		 {$90^{\circ}$}
		{\True $45^{\circ}$} 
		{$30^{\circ}$} 
		{$60^{\circ}$} 
	\loigiai{ 
		Đường thẳng $d_{1}\colon 3x+y-1=0$  có vectơ pháp tuyến $\overrightarrow{n}_{1}=(3;1)$.\\
		Đường  thẳng $d_{2}\colon 4x-2y-4=0$ có vectơ pháp tuyến $\overrightarrow{n}_{2}=(4;-2)$.\\
		 \[\cos \left(d_{1},d_{2}\right)=\left|\cos \left(\overrightarrow{n}_{1},\overrightarrow{n}_{2}\right)\right|=\dfrac{\left|\overrightarrow{n}_{1}\cdot \overrightarrow{n}_{2}\right|}{\left|\overrightarrow{n}_{1}\right|\cdot\left|\overrightarrow{n}_{2}\right|}=\dfrac{1}{\sqrt{2}}.\]
		Vậy  góc giữa hai đường thẳng $d_{1}$, $d_{2}$ là $45^{\circ}$.
	 }
\end{ex}

% nguồn: MaHoaNH_TracNghiem_Lop10
\begin{ex}%[0H9H3-4]%[ma:febcaf33]
	Trong mặt phẳng $Oxy$,  tính góc giữa hai đường thẳng $d\colon x-2y-1=0$ và $d'\colon x+3y-11=0$.
	\choice 
		{$30^\circ$} 
		 {\True $45^\circ$}
		 {$60^\circ$}
		 {$135^\circ$}
	 \loigiai{
		Đường thẳng $d$ có một véc-tơ pháp tuyến  là $\overrightarrow{n}=(1;-2)$.\\
		Đường thẳng $d'$ có một véc-tơ pháp tuyến là  $\overrightarrow{u}=(1;3)$.\\
		 Ta có $\cos\left(d,d'\right)=\dfrac{|1-6|}{\sqrt{1+4}\cdot \sqrt{1+9}}=\dfrac{5}{5\sqrt{2}}=\dfrac{\sqrt{2}}{2}$.\\
		Vậy góc giữa hai đường thẳng $d$ và $d'$ là  $45^\circ$.
	} 
\end{ex}

% nguồn: MaHoaNH_TracNghiem_Lop10
\begin{ex}%[0H9H3-4]%[ma:819761a0]
	Tính góc giữa hai đường  thẳng $d_{1}\colon x-3y+1=0$ và $d_{2}\colon x+2y-5=0$.
	 \choice
		{$60^{\circ}$} 
		{\True $45^{\circ}$} 
		{$135^{\circ}$} 
		 {$120^{\circ}$}
	\loigiai{ 
		Đường thẳng $(d_{1})$ có vectơ pháp  tuyến $\overrightarrow{n}_{1}=(1;-3)$.\\
		Đường thẳng $(d_{2})$ có vectơ pháp tuyến $\overrightarrow{n}_{2}=(1;  2)$.\\
		Ta có $\cos  \left(d_{1}, d_{2}\right)
		 =\left|\cos \left(\overrightarrow{n}_{1}, \overrightarrow{n}_{2}\right)\right|
		 =\dfrac{\left|\overrightarrow{n}_{1} \cdot \overrightarrow{n}_{2}\right|}{\left|\overrightarrow{n}_{1}\right| \cdot \left|\overrightarrow{n}_{2}\right|}
		=\dfrac{\left|1-6\right|}{\sqrt{10}  \cdot \sqrt{5}}
		=\dfrac{\sqrt{2}}{2}  \Rightarrow\left(d_{1}, d_{2}\right)=45^{\circ}$.
	 }
\end{ex}

% nguồn: MaHoaNH_TracNghiem_Lop10
\begin{ex}%[0H9H3-4]%[ma:f6b91124]
	Cho hai đường  thẳng $\left(d_1\right)\colon\heva{
	&x=2+t \\ 
	&y=5-2t}$ và $\left(d_2\right)\colon\heva{ 
	&x=4-s  \\
	&y=3-3s}$ ($t$, $s$  là các tham số). Góc giữa hai đường thẳng $d_1$ và $d_2$ bằng
	 \choice
		 {$90^{\circ}$}
		{\True $45^{\circ}$} 
		{$60^{\circ}$} 
		{$30^{\circ}$} 
	 \loigiai{
		Vectơ chỉ phương của $d_1$ là $\overrightarrow{u}_1=(1;-2)$, vectơ chỉ phương của  $d_2$ là $\overrightarrow{u}_2=(-1;-3)$.\\
		Gọi  $\varphi$ là góc giữa hai đường thẳng $d_1$ và $d_2$.\\
		Ta  có $\cos \varphi=\dfrac{\left|\overrightarrow{u}_1 \cdot \overrightarrow{u}_2\right|}{\left|\overrightarrow{u}_1\right| \cdot\left|\overrightarrow{u}_2\right|}=\dfrac{|1 \cdot(-1)+(-2) \cdot(-3)|}{\sqrt{1^2+(-2)^2} \cdot \sqrt{(-1)^2+(-3)^2}}=\dfrac{\sqrt{2}}{2}$.\\
		Do đó góc giữa hai đường thẳng $d_1$ và $d_2$ là $\varphi=45^{\circ}$. 
	 }
\end{ex}

% nguồn: MaHoaNH_TracNghiem_Lop10
\begin{ex}%[0H9Y3-4]%[ma:b9da320d]
	Người ta quy ước góc giữa hai đường  thẳng song song hoặc trùng nhau là
	 \choice
		 {$180^\circ$}
		 {$120^\circ$}
		 {$90^\circ$}
		{\True  $0^\circ$}
	\loigiai{ 
		Góc giữa hai  đường thẳng song song hoặc trùng nhau là $0^\circ$.
	 }
\end{ex}


%%%%%%%%%%%%%%%% ĐÚNG/SAI %%%%%%%%%%%%%%%%

% nguồn: MaHoaNH_DungSai_Lop10
\begin{ex}%[0H9N3-4]%[ma:00047697]
	Trong hệ tọa độ $Oxy$, giả sử $\alpha$ là góc giữa hai đường thẳng $d\colon x-\sqrt{3}y+2=0$ và $d'\colon x+\sqrt{3}y-1=0$. Các khẳng định sau đây  đúng hay sai?
	\choiceTF 
		 {\True $\sin\alpha=\dfrac{\sqrt{3}}{2}$}
		 {$\cos\alpha=-\dfrac{1}{2}$}
		{$\tan\alpha=-\sqrt{3}$} 
		{\True  $\cot\alpha=\dfrac{\sqrt{3}}{3}$}
	 \loigiai{
		Đường thẳng $d$, $d'$ lần lượt có véc-tơ pháp tuyến là $\overrightarrow{n}=(1;-\sqrt{3})$, $\overrightarrow{n'}=(1;\sqrt{3})$.\\ 
		 Khi đó $\cos\alpha=\dfrac{\left|1-3\right|}{\sqrt{1+3}\cdot\sqrt{1+3}}=\dfrac{1}{2}$ $\Rightarrow \alpha=60^\circ$.\\
		Do đó $\sin\alpha=\dfrac{\sqrt{3}}{2}$, $\cos\alpha=\dfrac{1}{2}$, $\tan\alpha=\sqrt{3}$ và  $\cot\alpha=\dfrac{\sqrt{3}}{3}$.
	} 
\end{ex}

% nguồn: MaHoaNH_DungSai_Lop10
\begin{ex}%[0H9H3-4]%[ma:45e19381]
	Trong mặt phẳng $Oxy$, cho hai  đường thẳng $d_1 \colon -x+2y-3=0$ và $d_2 \colon 2x-y+7=0$.
	\choiceTF 
		{Hai đường thẳng $d_1$ và $d_2$ vuông góc  nhau}
		{Tọa độ giao điểm của $d_1$ và $d_2$  là $\left(\dfrac{11}{3};\dfrac{1}{3}\right)$}
		{\True Cho $A(0;-4)$. Khoảng cách từ $A$  đến $d_1$ bằng khoảng cách từ $A$ đến $d_2$}
		 {\True Gọi $\varphi$ là góc giữa hai đường thẳng $d_1$ và $d_2$ thì $\cos \varphi = \dfrac{4}{5}$}
	\loigiai{ 
		\begin{itemchoice}
			\itemch Sai. Vì xét $d_1 \colon -x+2y-3=0$ có véc-tơ pháp tuyến là $\overrightarrow{n}_1=(-1;2)$ và $d_2 \colon 2x-y+7=0$ có véc-tơ pháp tuyến là $\overrightarrow{n}_2=(2;-1)$. Ta có $\overrightarrow{n}_1\cdot \overrightarrow{n}_2=(-1)\cdot 2+2\cdot (-1)=-4\neq 0$ suy  ra $d_1$ và $d_2$ không vuông góc nhau.
			\itemch Sai. Vì tọa độ giao điểm của $d_1$ và $d_2$ là nghiệm  của hệ $\heva{&-x+2y-3=0\\&2x-y+7=0}\Leftrightarrow\heva{&x=-\dfrac{11}{3}\\&y=-\dfrac{1}{3}.}$
			\itemch Đúng. Vì $\mathrm{d}(A,d_1)=\dfrac{|-0+2\cdot (-4)-3|}{\sqrt{(-1)^2+2^2}}=\dfrac{11}{\sqrt{5}}$,  $\mathrm{d}(A,d_2)=\dfrac{|2\cdot 0- (-4)+7|}{\sqrt{2^2+(-1)^2}}=\dfrac{11}{\sqrt{5}}$.
			 \itemch Đúng. Vì $\cos \varphi=\dfrac{|-1\cdot 1+2\cdot (-1)|}{\sqrt{(-1)^2+2^2}\cdot \sqrt{1^2+(-2)^2}}=\dfrac{4}{5}$.
		\end{itemchoice}
	} 
\end{ex}

% nguồn: MaHoaNH_DungSai_Lop10
\begin{ex}%[0H9H3-4]%[0H9H3-5]%[0H9N3-2]%[ma:3673d7e0]
	Trong mặt phẳng tọa độ $Oxy$, cho hai đường thẳng $\Delta_1\colon 2 x+y-1=0$ và $\Delta_2\colon \heva{&x=2+t \\ &y=1-t}$. Xét tính đúng sai của mỗi  khẳng định sau:
	\choiceTF 
		{Đường thẳng $\Delta_2$ có một vec-tơ chỉ phương là $\vec{u}_2=(2;  1)$}
		{Đường thẳng $\Delta_1$ có một vec-tơ pháp tuyến là $\vec{n}_1=(2;  1)$ nên $\Delta_1$ có một vec-tơ chỉ phương là $\vec{u}_1=(1; 2)$}
		{\True Khoảng cách từ điểm $M(2; 1)$ đến đường thẳng  $\Delta_1$ bằng $\dfrac{4}{\sqrt{5}}$}
		{\True Cosin của góc tạo bởi hai  đường thẳng $\Delta_1$ và $\Delta_2$ bằng $\dfrac{3}{\sqrt{10}}$}
	 \loigiai{
		\begin{itemchoice}
			\itemch \textbf{Sai}.  Vì đường thẳng $\Delta_2$ nhận $\vec{u}_2=(1;-1)$ làm một vec-tơ chỉ phương.
			\itemch \textbf{Sai}. Vì đường thẳng $\Delta_1$ nhận $\vec{u}_1=(-1;2)$ làm một vec-tơ chỉ  phương.
			\itemch \textbf{Đúng}. Vì $\mathrm{d}\big(  M, \Delta_1\big)=\dfrac{|2\cdot (2)+(1)-1|}{\sqrt{2^2+1^2}}=\dfrac{4}{\sqrt{5}}$.
			\itemch \textbf{Đúng}. Vì $\cos  \left(\Delta_1, \Delta_2 \right)=\dfrac{|-1\cdot1+2\cdot(-1)|}{\sqrt{2^2+1^2}\cdot\sqrt{1^2+1^2}}=\dfrac{3}{\sqrt{10}}$.
		\end{itemchoice}
	} 
\end{ex}

% nguồn: MaHoaNH_DungSai_Lop10
\begin{ex}%[0H9H3-4]%[ma:86ef87c0]
	Cho hai  đường thẳng $\Delta_1\colon x-2y+15=0$ và $\Delta_2\colon\heva{&x=2-t\\&y=4+2t}(t\in\mathbb{R})$. Xét tính đúng - sai của các khẳng định bên dưới.
	 \choiceTF
		{Đường thẳng $\Delta_1$ cắt trục tung tại  điểm $A(0;-15)$}
		{\True Đường thẳng  $\Delta_2$ đi qua điểm $M(1;6)$}
		{Hai đường thẳng $\Delta_1$ và $\Delta_2$ song song nhau} 
		{\True Góc giữa $\Delta_1$ và $\Delta_2$ là  $90^\circ$}
	\loigiai{$\Delta_1$ có véc-tơ pháp tuyến $\overrightarrow{n}_{1}=(1;-2)$.\\ 
		$\Delta_2$ có véc-tơ chỉ phương $\overrightarrow{u}_{2}=(-1;2)$ nên có  véc-tơ pháp tuyến $\overrightarrow{n}_{2}=(2;1)$.
		\begin{itemchoice}
			\itemch Khi $x=0$, thay vào $\Delta_1$ ta có $y=\dfrac{15}{2}$ nên $\Delta_1$ cắt trục tung tại điểm $\left(0;\dfrac{15}{2}\right)$, suy ra đường thẳng $\Delta_1$ cắt trục tung tại điểm $A(0;-15)$ là  sai.
			\itemch Thay tọa độ $M(1;6)$ vào phương trình $\Delta_2$ ta  có $\heva{&1=2-t\\&6=4+2t}\Leftrightarrow\heva{&t=1\\&t=1}$, vậy suy ra đường thẳng $\Delta_2$ đi qua điểm $M(1;6)$ là đúng.
			\itemch Ta có $\overrightarrow{n}_1$ không cùng phương $\overrightarrow{n}_2$ nên hai đường  thẳng $\Delta_1$ và $\Delta_2$ song song nhau là sai.
			\itemch Ta có $\overrightarrow{n}_1\cdot\overrightarrow{n}_2=0$, suy ra $\Delta_1\perp\Delta_2$  nên góc giữa $\Delta_1$ và $\Delta_2$ là $90^\circ$ là đúng.
	\end{itemchoice}}
\end{ex}

% nguồn: MaHoaNH_DungSai_Lop10
\begin{ex}%[0H9H3-4]%[ma:5585c17a]
	Trong hệ tọa độ $Oxy$, cho các đường thẳng $d_1\colon 2x-y-10=0$, $d_2\colon x-3y+9=0$ và $d_3\colon x+2y-1=0$. Gọi $\alpha$, $\beta$ lần lượt là góc  giữa hai đường thẳng $d_1$, $d_2$ và $d_1$, $d_3$. Các khẳng định sau đây đúng hay sai?
	\choiceTF 
		{$\alpha>\beta$} 
		{\True $\beta=2\alpha$} 
		 {\True $3\alpha-\beta=45^\circ$}
		 {$\alpha+\beta=90^\circ$}
	 \loigiai{
		Các đường thẳng $d_1$, $d_2$, $d_3$ lần lượt có véc-tơ pháp tuyến là $\overrightarrow{n_1}=(2;-1)$, $\overrightarrow{n_2}=(1;-3)$, $\overrightarrow{n_3}=(1;2)$.\\ 
		Khi  đó
		\begin{itemize}
			 \item $\cos\alpha=\dfrac{\left|2\cdot 1+(-1)\cdot (-3)\right|}{\sqrt{2^2+(-1)^2}\cdot\sqrt{1^2+(-3)^2}}=\dfrac{1}{\sqrt{2}}$ $\Rightarrow \alpha=45^\circ$.
			\item  $\cos\beta=\dfrac{\left|2\cdot 1+(-1)\cdot 2\right|}{\sqrt{2^2+(-1)^2}\cdot\sqrt{1^2+2^2}}=0$ $\Rightarrow \beta=90^\circ$.
		\end{itemize}
		Từ các phương án ta thấy các khẳng định đúng là $\beta=2\alpha$, $3\alpha-\beta=45^\circ$. 
	 }
\end{ex}


%%%%%%%%%%%%%%%% TRẢ LỜI NGẮN %%%%%%%%%%%%%%%%

% nguồn: MaHoaNH_TraLoiNgan_Lop10
\begin{bt}%[0H9H3-4]%[ma:b577a8d1]
	Trên mặt phẳng  tọa độ $Oxy$, cho hai đường thẳng $d_1\colon x-2y+3=0$ và $d_2\colon (m+1)x+(m-1)y+2=0$ với $m$ là tham số. Tìm giá trị của $m$ để hai đường thẳng vuông góc với nhau.
	\shortans{3} 
	 \loigiai{
		Ta có véc-tơ pháp tuyến của $d_1$ là $\overrightarrow{n}_1=(1;-2)$; véc-tơ pháp tuyến của $d_2$ là  $\overrightarrow{n}_2=(m+1;m-1)$.\\
		$d_1\perp d_2\Leftrightarrow\overrightarrow{n}_1\cdot\overrightarrow{n}_2=0\Leftrightarrow1(m+1)-2(m-1)=0\Leftrightarrow-m+3=0\Leftrightarrow m=3$. 
	 }
\end{bt}

% nguồn: MaHoaNH_TraLoiNgan_Lop10
\begin{bt}%[0H9H3-4]%[ma:a7c75f18]
	Trong mặt phẳng $Oxy$, tìm giá trị  của tham số $m$ để hai đường thẳng $d_1\colon (2 m-1) x+m y-10=0$ và $d_2\colon x+2 y+6=0$ vuông góc nhau.
	\shortans{0,25}  \loigiai{
		Đường  thẳng $d_1\colon (2 m-1) x+m y-10=0$ có một vec-tơ pháp tuyến $\vec{n}_1=(2 m-1; m)$.\\
		Đường thẳng $d_2\colon x+2 y+6=0$ có một vec-tơ  pháp tuyến $\vec{n}_2=(1; 2)$.\\
		Hai đường thẳng $d_1 \perp  d_2 \Leftrightarrow \vec{n}_1 \cdot \vec{n}_2=0 \Leftrightarrow(2 m-1)+2 m=0 \Leftrightarrow m=\dfrac{1}{4}=0{,}25$.
	} 
\end{bt}

% nguồn: MaHoaNH_TraLoiNgan_Lop10
\begin{bt}%[0H9V3-4]%[ma:bd2c0d74]
	Trong mặt phẳng $Oxy$, cho hai đường thẳng $d\colon x-\sqrt{3} y+2=0$ và $\Delta\colon \heva{&x=1+t \\ &y=-2+m t}$ $(t \in \mathbb{R})$ với $m$ là tham số thực. Giá trị của  $m$ để góc giữa hai đường thẳng $d$ và $\Delta$ bằng $60^{\circ}$ là bao nhiêu? (Kết quả làm tròn đến một chữ số sau dấu phẩy).
	 \shortans{-0,6}
	\loigiai{ 
		Ta có $\overrightarrow{n}_d=\left( 1;-\sqrt{3}\right) \Rightarrow  \overrightarrow{u}_d=\left( \sqrt{3}; 1\right) $ và $\overrightarrow{u}_{\Delta}=(1; m)$.
		 Theo đề bài
		\begin{eqnarray*}
			\cos (d, \Delta)=\dfrac{\left| \sqrt{3}+m\right| }{2 \cdot \sqrt{1+m^2}} & \Leftrightarrow& \cos 60^{\circ}=\dfrac{1}{2} \cdot \dfrac{\left| m+\sqrt{3}\right| }{\sqrt{m^2+1}} \Leftrightarrow \sqrt{m^2+1}=\left| m+\sqrt{3}\right| \\ 
			&  \Leftrightarrow& m=-\dfrac{\sqrt{3}}{3}.
		\end{eqnarray*}
		Vậy $m=-\dfrac{\sqrt{3}}{3}  \approx -0{,}6$ thỏa mãn đề bài.
	} 
\end{bt}
